\chapter*{Resumen}
\makeatletter\@mkboth{Resumen}{Resumen}\makeatother
\addcontentsline{toc}{chapter}{Resumen}

En la última década, el aprendizaje computacional se ha consolidado como el motor fundamental detrás de avances revolucionarios en áreas como la visión artificial, el procesamiento del lenguaje natural y la robótica autónoma. No obstante, el éxito de los modelos más potentes, particularmente en el paradigma del aprendizaje supervisado, depende críticamente de la disponibilidad de vastos conjuntos de datos etiquetados. El proceso de etiquetado, a menudo realizado por expertos humanos, constituye un cuello de botella logístico y económico, siendo en muchos casos una tarea laboriosa, lenta y propensa a errores. Ante este escenario, el aprendizaje semi-supervisado (SSL, por sus siglas en inglés) se presenta como una disciplina esencial que busca maximizar el rendimiento de los clasificadores utilizando un conjunto reducido de datos etiquetados en combinación con una gran cantidad de datos no etiquetados.

Este trabajo aborda el problema de la clasificación semi-supervisada desde una perspectiva analítica, analizando cómo el conocimiento de la geometría del espacio de entrada permite refinar las fronteras de decisión. A diferencia de otros enfoques más empíricos, este estudio se centra en las bases matemáticas que justifican el uso de datos sin etiquetas. Para que estos datos aporten valor, deben cumplirse ciertas hipótesis de regularidad. En este sentido, se formalizan y discuten las tres suposiciones fundamentales del SSL: la suposición de suavidad, que postula que puntos cercanos en regiones de alta densidad deben compartir la misma etiqueta; la suposición de clúster o de baja densidad, que sugiere que la frontera de decisión debe pasar por regiones donde los datos son escasos; y la suposición de variedad, que asume que los datos de alta dimensión se concentran en una estructura de menor dimensión. Comprender estas hipótesis es crucial, ya que si los datos no las satisfacen, la inclusión de ejemplos no etiquetados puede introducir ruido y degradar el rendimiento del modelo original.

La contribución central de este trabajo es la propuesta de una taxonomía estructurada que organiza la vasta literatura del aprendizaje semi-supervisado, inspirada en revisiones recientes (en particular, van Engelen y Hoos, 2020) y adaptada con modificaciones motivadas a los objetivos del estudio. A través de un enfoque didáctico pero riguroso, se clasifican los algoritmos atendiendo a su naturaleza matemática y a la forma en que incorporan la información no supervisada. El trabajo no solo recopila información dispersa, sino que la sintetiza mediante un análisis comparativo de las funciones de pérdida y los términos de regularización específicos de cada paradigma. El objetivo es proporcionar al lector un mapa conceptual que permita discernir entre métodos inductivos, diseñados para aprender una regla de decisión generalizable a nuevos datos, y métodos transductivos, cuyo propósito es inferir etiquetas exclusivamente para un conjunto de puntos no etiquetados previamente definidos.

La taxonomía propuesta se desarrolla detalladamente a través de cuatro grandes bloques. En primer lugar, se analizan los métodos envolventes, como el auto-entrenamiento y los métodos basados en desacuerdo (co-regularización y múltiples vistas), que utilizan modelos supervisados de forma iterativa para generar pseudo-etiquetas. En segundo lugar, se estudia el preprocesado no supervisado, donde se emplean técnicas de extracción de características y clustering para transformar el espacio de representación antes de aplicar el clasificador. El tercer bloque, de especial interés matemático, analiza los métodos intrínsecamente semi-supervisados. Aquí se profundiza en modelos generativos que utilizan el algoritmo EM para optimizar mezclas de distribuciones, y en métodos discriminativos como las Máquinas de Vectores de Soporte Semi-supervisadas (S3VM), que buscan maximizar el margen en presencia de datos no etiquetados. Finalmente, se exploran los métodos basados en grafos, analizando la construcción de matrices de adyacencia y el uso de funciones armónicas y cortes mínimos (mincut) para la propagación de etiquetas.

En conclusión, este trabajo establece un marco teórico sólido que clarifica los fundamentos matemáticos de las técnicas actuales de aprendizaje semi-supervisado. La revisión sistemática de los algoritmos y sus propiedades pone de manifiesto que el éxito del SSL no depende de la cantidad de datos per se, sino de la correcta alineación entre las suposiciones del modelo y la estructura geométrica subyacente de los datos. Este estudio sienta las bases para futuras líneas de investigación, como la evaluación empírica de estas taxonomías en grandes modelos de lenguaje o sistemas de visión avanzados, garantizando una transición fluida entre la teoría matemática y su implementación práctica.